% header.tex — LaTeX preamble additions injected by Pandoc (--include-in-header).
% These run after Pandoc's own \usepackage lines, so packages it already loads
% (longtable, booktabs, graphicx, hyperref, xcolor) are available here.

% --- Code blocks / ASCII diagrams (wrap-code.lua emits fvextra Verbatim) -----------
\usepackage{fvextra}            % loads fancyvrb; adds breaklines/breakanywhere
% breakanywhere is NOT a default: wrap-code.lua adds it per block, only where a line has an
% unbreakable run wider than the page (it exhausts TeX's main memory if used on every block).
\fvset{fontsize=\footnotesize,breaklines=true,breaksymbolleft={}}

% --- Tables ------------------------------------------------------------------------
\usepackage{etoolbox}
\usepackage{ragged2e}
% Pandoc lays out wrapping table columns with \raggedright, which disables
% hyphenation, so long compound words ("Spectre/retpoline", "non-speculative")
% overflow narrow cells. Make \raggedright behave like \RaggedRight, which keeps
% hyphenation, so cell text breaks to fit.
\let\raggedright\RaggedRight
% Shrink wide tables, trim inter-column padding, and allow last-resort stretch so
% cells fit the text block. The font/padding are redefinable macros so the
% fit-tables.lua filter can drop especially wide tables (many columns) to a smaller
% size and tighter spacing without affecting the normal tables.
\newcommand{\tablefontsize}{\footnotesize}
\newcommand{\tablecolsep}{4pt}
\AtBeginEnvironment{longtable}{\tablefontsize\setlength{\tabcolsep}{\tablecolsep}\setlength{\emergencystretch}{2em}}
\setlength{\LTpre}{0.6em}
\setlength{\LTpost}{0.6em}
\renewcommand{\arraystretch}{1.15}

% --- Blank verso pages truly blank (no header/footer on openright blanks) -----------
\usepackage{emptypage}

% --- Typography --------------------------------------------------------------------
\usepackage{microtype}
\usepackage{seqsplit}           % allow ugly-long inline tokens to break
\sloppy                         % prefer extra interword space over overfull lines

% Map the technical/math Unicode glyphs used in the prose onto symbols the body font
% (Latin Modern) has, so they render instead of silently dropping out. These mappings
% are inactive inside Verbatim, so the box-drawing diagrams in code keep their literal
% Menlo glyphs. (§, accented letters, dashes, quotes, ×, … already render in LM.)
\usepackage{amssymb}
\usepackage{textcomp}
\usepackage{pifont}
\usepackage{newunicodechar}
% relations / operators
\newunicodechar{≈}{\ensuremath{\approx}}
\newunicodechar{≠}{\ensuremath{\neq}}
\newunicodechar{≤}{\ensuremath{\leq}}
\newunicodechar{≥}{\ensuremath{\geq}}
\newunicodechar{≪}{\ensuremath{\ll}}
\newunicodechar{≫}{\ensuremath{\gg}}
\newunicodechar{±}{\ensuremath{\pm}}
\newunicodechar{÷}{\ensuremath{\div}}
\newunicodechar{−}{\ensuremath{-}}
\newunicodechar{∈}{\ensuremath{\in}}
% arrows
\newunicodechar{→}{\ensuremath{\rightarrow}}
\newunicodechar{←}{\ensuremath{\leftarrow}}
\newunicodechar{↔}{\ensuremath{\leftrightarrow}}
\newunicodechar{↓}{\ensuremath{\downarrow}}
\newunicodechar{↑}{\ensuremath{\uparrow}}
\newunicodechar{⇒}{\ensuremath{\Rightarrow}}
% marks and stars (tables / bibliography)
\newunicodechar{★}{\ding{72}}
\newunicodechar{✓}{\ding{51}}
\newunicodechar{✗}{\ding{55}}
\newunicodechar{®}{\textregistered}
% greek used in prose
\newunicodechar{α}{\ensuremath{\alpha}}
\newunicodechar{κ}{\ensuremath{\kappa}}
\newunicodechar{Σ}{\ensuremath{\Sigma}}
\newunicodechar{μ}{\textmu}          % U+03BC Greek mu, used as "μs" (micro); distinct from µ U+00B5
% super/subscripts
\newunicodechar{⁰}{\textsuperscript{0}}
\newunicodechar{³}{\textsuperscript{3}}
\newunicodechar{⁴}{\textsuperscript{4}}
\newunicodechar{⁵}{\textsuperscript{5}}
\newunicodechar{⁶}{\textsuperscript{6}}
\newunicodechar{⁸}{\textsuperscript{8}}
\newunicodechar{⁻}{\textsuperscript{\ensuremath{-}}}
\newunicodechar{₂}{\textsubscript{2}}
% glyphs the chapters added in the 2026-11 insertion rounds use and Latin Modern lacks
\newunicodechar{⁷}{\textsuperscript{7}}
\newunicodechar{⁹}{\textsuperscript{9}}
\newunicodechar{ⁿ}{\textsuperscript{n}}
\newunicodechar{ᵇ}{\textsuperscript{b}}
\newunicodechar{₀}{\textsubscript{0}}
\newunicodechar{₁}{\textsubscript{1}}
\newunicodechar{φ}{\ensuremath{\varphi}}
\newunicodechar{λ}{\ensuremath{\lambda}}
\newunicodechar{σ}{\ensuremath{\sigma}}
\newunicodechar{τ}{\ensuremath{\tau}}
\newunicodechar{∪}{\ensuremath{\cup}}
\newunicodechar{⊇}{\ensuremath{\supseteq}}
\newunicodechar{◆}{\ding{117}}

% --- Running heads -----------------------------------------------------------------
\usepackage{fancyhdr}
\setlength{\headheight}{14pt}
\pagestyle{fancy}
\fancyhf{}
\renewcommand{\chaptermark}[1]{\markboth{\thechapter.\ #1}{}}
\renewcommand{\sectionmark}[1]{\markright{\thesection.\ #1}}
\fancyhead[LE]{\small\nouppercase{\leftmark}}    % verso: chapter
\fancyhead[RO]{\small\nouppercase{\rightmark}}   % recto: section
\fancyhead[RE,LO]{}
\fancyfoot[LE,RO]{\small\thepage}
\renewcommand{\headrulewidth}{0.4pt}
\renewcommand{\footrulewidth}{0pt}
% Chapter/part opening pages use plain style; keep the page number there:
\fancypagestyle{plain}{%
  \fancyhf{}%
  \fancyfoot[LE,RO]{\small\thepage}%
  \renewcommand{\headrulewidth}{0pt}%
}

% --- Table of contents: room for three-digit chapter numbers ---------------------------
% The book class sets a chapter number in a 1.5em box and a section number in 2.3em: enough
% for "99" and "9.9", not for "100" and "100.4" (the number ran into the title once the book
% passed 99 chapters). Widen both.
\makeatletter
\patchcmd{\l@chapter}{\setlength\@tempdima{1.5em}}{\setlength\@tempdima{2.4em}}{}{}
\renewcommand*\l@section{\@dottedtocline{1}{2.4em}{3.1em}}
\makeatother

% --- Tighten Pandoc lists a touch --------------------------------------------------
\providecommand{\tightlist}{\setlength{\itemsep}{0pt}\setlength{\parskip}{0pt}}
